% Source: https://artofproblemsolving.com/community/c3h2332487_amc_8_prep_4_mustknow_formulas
% Stable identity: aops-topic-2332487-q01
\textbf{(a)} By the symmetry $\binom Nr=\binom N{N-r}$ with $N=n+k-1$ and $r=k$, we get $\binom{n+k-1}{k}=\binom{n+k-1}{n-1}$, so the two expressions agree. For $n=3$ and $k=2$ they give $\binom42=6$. But the ways to split $3$ identical objects between $2$ bins with none empty are only $(1,2)$ and $(2,1)$: there are $2$ ways, not $6$. So the statement is false.

\textbf{(b)} Line up the $n$ objects in a row. There are $n-1$ gaps between neighbouring objects. Choosing $k-1$ of these gaps and putting a divider in each splits the row into $k$ nonempty blocks, which we put into bins $1,2,\dots,k$ in order. Conversely, every placement with all bins nonempty arises from exactly one choice of gaps (the dividers go after the objects of bin $1$, after those of bin $2$, and so on). Hence the number of ways is
\[\boxed{\binom{n-1}{k-1}}.\]
(For $n=3$, $k=2$ this is $\binom21=2$, as found in (a).)

\textbf{(c)} Placements of $n$ objects into $k$ bins, empty bins allowed, correspond one-to-one to placements of $n+k$ objects with every bin nonempty: add one extra object to each bin, or remove one from each. By (b) their number is
\[\binom{n+k-1}{k-1}=\binom{n+k-1}{n}.\]
So the guide's upper index $n+k-1$ belongs to the count \emph{with} empty bins allowed, and even there the lower index must be $k-1$ (or $n$), not $k$. For example, $3$ objects in $2$ bins, empty allowed, give $\binom41=4$ ways: $(0,3),(1,2),(2,1),(3,0)$.
